\chapter{Related Work}
\label{chap:related}

\section{Reasoning with a single language model}

Modern language-model reasoning systems often begin with chain-of-thought prompting, which elicits intermediate natural-language steps before an answer \citep{wei2022chain}. The important methodological insight is that reasoning quality depends not only on the final string but on the trajectory that produces it. Self-consistency extends this idea by sampling multiple reasoning paths and selecting the most consistent answer, thereby separating candidate generation from aggregation without requiring persistent agent identities \citep{wang2023selfconsistency}. Iterative approaches such as Self-Refine and Reflexion add feedback and revision loops to a single model \citep{madaan2023selfrefine,shinn2023reflexion}. These methods demonstrate that repeated computation can help, but they also expose a general risk shared with multi-agent systems: later text can overwrite a correct earlier state.

This dissertation uses a single-agent system (SAS) as the primary baseline because architecture comparisons are only meaningful relative to a strong, compute-aware alternative. A weak one-shot baseline would confound decomposition with additional inference. The SAS receives the same question and oracle context as the multi-agent systems and is evaluated at the same nominal thinking-token levels. This follows the matched-compute emphasis of \citet{tran2026single}, who argue that the apparent benefit of multi-agent reasoning can diminish or reverse once the comparison controls for thinking tokens.

Single-agent self-consistency is related to the Ensemble architecture studied here, but the decision rules differ. Self-consistency normally aggregates answers by agreement, while this implementation exposes several candidate generations to a final language-model aggregator. The latter can in principle use evidence quality rather than frequency, but it can also introduce selection error. This distinction motivates the oracle-final gap used later: the gap measures how often an answer was available among candidates but was not converted into a correct final prediction.

\section{Multi-agent language-model systems}

Multi-agent LLM systems distribute inference across interacting calls that are described as agents. The term covers substantially different computational graphs. Some systems use independent roles to diversify viewpoints, some use planners and workers to decompose tasks, and others use adversarial critique. AutoGen provides a general conversation-based framework in which specialised agents exchange messages and invoke tools \citep{wu2023autogen}. Such flexibility is useful for applications, but it makes scientific comparison difficult: changing the number of calls, prompts, stopping conditions, and aggregation logic can change both compute and behaviour.

Multi-agent debate is one prominent family. \citet{du2023debate} report that multiple model instances proposing and debating answers can improve factuality and reasoning. \citet{liang2023mad} similarly motivate debate as a mechanism for divergent thinking and mitigation of degenerate consensus. At a high level, debate has two expected benefits. First, independent proposals increase the probability that at least one agent finds a correct answer. Second, revision after observing another proposal may correct mistakes. However, debate also adds two decision points: revisions may turn correct proposals into wrong ones, and a final picker may choose the wrong revised answer. The present study measures both points rather than treating debate as a single black box.

Planner-worker architectures make a different promise. A planner decomposes a multi-hop question into subquestions; workers solve these steps either sequentially or in parallel; and an aggregator forms the final answer. Sequential execution can propagate useful evidence from one worker to the next, but later workers inherit earlier errors and may displace previously correct content. Parallel execution reduces error propagation and latency but deprives workers of intermediate findings. Role-based parallelism seeks diversity by asking several agents to analyse the same evidence from different perspectives. An ensemble samples multiple answers with stochastic decoding. These architectures should not be treated as interchangeable instances of ``more agents'' because they create different source states and different opportunities for loss.

The dissertation therefore treats architecture as a temporal data-generating process. Debate is proposal $\rightarrow$ revision $\rightarrow$ picker. Sequential is plan $\rightarrow$ worker chain $\rightarrow$ aggregator. Parallel Roles and Ensemble are candidate set $\rightarrow$ aggregator. Subtask Parallel is plan $\rightarrow$ partial findings $\rightarrow$ aggregator. The source-stage semantics differ: a subtask finding is not necessarily expected to answer the full question, whereas a revised debate answer is. This is why the strongest human transition comparison is restricted to Debate and Sequential, where the source and target states are interpretable and the final boundary is explicit.

\section{Multi-hop question answering}

Multi-hop QA is a natural test bed for decomposition because the answer is intended to require combining facts. MuSiQue was constructed by composing single-hop questions while applying filters designed to reduce shortcut solutions \citep{trivedi2022musique}. 2WikiMultiHopQA provides questions requiring reasoning over evidence from multiple Wikipedia pages and includes supporting facts intended to expose reasoning steps \citep{ho2020twowiki}. The datasets differ in difficulty and construction, which makes their joint use valuable: an architectural pattern that appears on only one dataset may reflect dataset-specific answer forms or shortcut opportunities.

At the same time, benchmark labels do not prove that a model used the intended reasoning chain. \citet{min2019compositional} showed that compositional questions can often be answered without the expected multi-hop process. For this reason, the present dissertation does not claim to verify latent reasoning. It evaluates observable text states: whether an answer-equivalent claim appears, whether the system endorses it, and whether later components preserve it. Oracle context is supplied so that the experiment focuses on reasoning and coordination rather than retrieval. This narrows the claim but improves internal interpretability: absence of the answer cannot be attributed to a failed document retriever.

Exact match and token-level F1 remain useful because they are transparent and reproducible, but they are brittle for semantically equivalent answers. An output may include the correct entity with explanatory text, an alias, a date format, or a harmless modifier and still fail strict matching. Conversely, substring matching can over-credit an answer mentioned only to reject it. These limitations are amplified at intermediate stages, whose outputs are not written in final-answer format. The project consequently reports lexical final accuracy for continuity with standard QA evaluation and uses a separate semantic protocol for stage localisation.

\section{Process evaluation and failure localisation}

Outcome supervision assigns a label to a completed response. Process supervision evaluates intermediate steps, potentially detecting where an otherwise promising trajectory goes wrong. Work such as \citet{lightman2023verify} demonstrates the value of step-level feedback for mathematical reasoning. The setting here is related but not identical. Intermediate stages are not individual mathematical derivations; they are messages produced by different architectural roles. A useful label must therefore be role-aware and must distinguish informational presence from the system's decision.

The dissertation introduces two observable variables for every relevant stage. Discovery records whether the reference answer is substantively available in the text. Endorsement records whether the text's answer-bearing statement commits to it. Consider a revised debate output stating that ``Paris is plausible, but the evidence instead supports Lyon.'' Paris has been mentioned but is not endorsed. A presence-only judge may call this correct even though the revision would direct a final picker towards Lyon. Conversely, a worker may provide the exact supporting relationship without formatting a final answer; this can count as discovery even when endorsement is not applicable or remains uncertain. Explicit decision rules are therefore necessary for reproducible transition labels.

The closest quantitative concept is oracle accuracy over a candidate set. If any candidate is correct, an oracle selector would succeed. The difference between oracle and final accuracy is often described as aggregation regret. This project extends the idea temporally. It asks not only whether a correct candidate exists, but whether it survives revision, worker propagation, and final selection. Harmful and beneficial transitions then measure the net behaviour of a coordination step, rather than assuming that aggregation is beneficial whenever candidate coverage is high.

\section{Language models as evaluators}

LLM-based evaluation can approximate human preferences and semantic judgements at a scale that manual annotation cannot match. MT-Bench and Chatbot Arena helped establish the practical use of strong models as judges while documenting biases such as position sensitivity and self-preference \citep{zheng2023judge}. G-Eval uses structured criteria and chain-of-thought evaluation to improve alignment with human judgements on generation tasks \citep{liu2023geval}. These results support machine-assisted evaluation, but they do not justify assuming validity on a new label definition or a deliberately difficult subpopulation.

Reliability is especially important for transition analysis. Suppose a source stage is falsely marked correct and the destination is correctly marked wrong. The resulting \ctow{} transition is entirely artificial. Two individually plausible labels can also form the wrong transition class when one error occurs. General calibration accuracy therefore does not guarantee transition-level accuracy. The present project first calibrated a Mistral judge on ordinary evidence-validity examples, then conducted a targeted human audit on transition-enriched strata. The targeted result was materially worse than the general result, so the dissertation does not use the machine judge as the primary estimator for Debate and Sequential transitions.

Cohen's kappa is reported alongside exact agreement because it adjusts for chance agreement under the observed marginal distributions \citep{cohen1960coefficient}. Bootstrap confidence intervals describe uncertainty in agreement estimates \citep{efron1994bootstrap}. Neither statistic establishes construct validity by itself. Hidden repeats measure one annotator's internal consistency, not independent inter-annotator agreement, and high consistency would not correct a systematically wrong codebook. The audit is therefore designed around explicit evidence checks, blind identifiers, frozen strata, and unblinding only after labels are complete.

\section{Positioning of this dissertation}

Existing work provides strong reasons to test both the promise and the cost of multi-agent decomposition. Debate and ensemble methods motivate diverse candidate generation; matched-compute studies question whether that diversity is worth its inference cost; process evaluation motivates inspection of intermediate states; and LLM-as-judge research offers scale with known validity risks. The gap is a controlled, architecture-aware analysis connecting these strands. This dissertation fills that gap by comparing multiple architectures under common budgets, measuring discovery and final correctness separately, localising loss at specific boundaries, and using human labels to support the claims for which machine errors would be most damaging.
